% !TeX program = lualatex
\documentclass[11pt]{article}
\usepackage[margin=25mm]{geometry}
\usepackage{fontspec}
\setmainfont{Latin Modern Roman}
\usepackage{amsmath,amssymb,mathtools}
\usepackage{unicode-math}
\setmathfont{Latin Modern Math}
\usepackage{xurl,hyperref}
\hypersetup{colorlinks=true,urlcolor=blue}
\setcounter{secnumdepth}{0}
\setlength{\parindent}{0pt}
\setlength{\parskip}{5pt}
\setlength{\emergencystretch}{3em}
\newcommand{\R}{\mathbb R}
\newcommand{\N}{\mathbb N}
\newcommand{\Z}{\mathbb Z}
\newcommand{\Q}{\mathbb Q}
\newcommand{\C}{\mathbb C}
\newcommand{\D}{\mathbb D}
\newcommand{\F}{\mathbb F}
\newcommand{\sref}[2]{\hyperlink{source:#1:#2}{[S#2]}}
\newcommand{\cataloguescope}[1]{#1}
\begin{document}
\section{Spectrahedrality of Hyperbolicity Cones}\label{statement}

\subsection{Definitions and mathematical statement}

\paragraph{Shared notation.}$\mathbb N$, $\mathbb Z$, $\mathbb Q$, $\mathbb R$, and $\mathbb C$ denote the natural numbers, integers, rationals, reals, and complex numbers, respectively. All additional parameters and hypotheses are specified in the question.


\paragraph{Statement recorded in the supplied source.}
Let $h$ be a real homogeneous polynomial hyperbolic in a direction $e$, meaning that $h(e)>0$ and $t\mapsto h(x-te)$ has only real zeros for every real vector $x$. Must its closed hyperbolicity cone be expressible as $\{x:\sum_i x_iA_i\succeq0\}$ for real symmetric matrices $A_i$? Such a cone is spectrahedral.
\par\smallskip\textit{Formulation references:} \sref{statement}{1}, \sref{statement}{2}.
\subsection{Short English statement}
Let $h$ be a real homogeneous polynomial hyperbolic in a direction $e$, meaning that $h(e)>0$ and $t\mapsto h(x-te)$ has only real zeros for every real vector $x$. Must its closed hyperbolicity cone be expressible as $\{x:\sum_i x_iA_i\succeq0\}$ for real symmetric matrices $A_i$? Such a cone is spectrahedral.
\subsection{Sources}
\begin{itemize}
\item[\hypertarget{source:statement:1}{S1}] ULAM.AI and the UnsolvedMath Contributors, UnsolvedMath: A Curated Collection of Open Mathematics Problems, catalogue entry. Catalogue attribution under CC BY 4.0.
\url{https://huggingface.co/datasets/ulamai/UnsolvedMath}.
\item[\hypertarget{source:statement:2}{S2}] Original problem formulation and mathematical context.
\url{https://doi.org/10.4171/owr/2014/17}.
\end{itemize}
\label{end:statement}
\end{document}
